\chapter{长期价值与连续决策}
\label{chap:rs-long-term}

用户正在挑选通勤背包，连续点击了几个相似款式。继续展示同类商品，可能仍有较高的点击概率；也可能让用户反复比较却迟迟不能作出决定。插入一个水壶或相机，又可能是无关干扰，也可能帮助系统判断用户究竟在准备日常通勤还是一次旅行。一次展示的价值不只来自当次响应，还来自它怎样改变后续需求、可用信息和选择机会。

前面的章节已经建立了用户表示、候选、响应预测、展示行动以及多方目标。本章增加的不是更长的输入序列，而是行动对未来的影响。问题定义要说清哪些状态会变，价值估计要回答一个行动在后续策略下值多少，策略优化要把估计或反馈变成可执行的选择，效果检验则要判断改善究竟由什么证据支持。这几方面共同决定了“长期”是否有明确含义，而不是一种模型名称。

\section{连续决策问题}
\label{sec:rs-sequential-problem}

把一次展示当作独立样本时，系统通常计算当前请求下的响应。\term{连续决策}（sequential decision making）则把多次行动连起来：本次选择既产生反馈，也改变下一次选择所面对的条件。这里的$t$表示决策时点，$a_t$沿用第~\ref{chap:rs-presentation}~章的展示行动，可以包含列表$L_t$、位置、发送时间或交互方式。一次决策可以对应一次刷新、一次会话或一天的分配；不同时间粒度不能在累计目标中不加区别地相加。

为描述这种联系，记$s_t$为环境在决策前的充分状态，即给定它和当前行动后，预测下一步不再需要更早的环境历史。状态可以包含当前需求、兴趣、重复曝光疲劳、提供者的生产能力和剩余资源。环境产生奖励$r_t$与下一状态$s_{t+1}$，其条件分布记为$P(r_t,s_{t+1}\mid s_t,a_t)$。可行行动范围$\mathcal{F}_t$由第~\ref{sec:rs-decision-resources}~节的资格与资源要求确定；如果预算剩余多少会改变未来选择，预算就必须随状态一起更新。

充分状态不等于系统实际能看到的信息。兴趣与疲劳往往不能直接观测，请求时可用的仍是$\bm{x}_t$和历史$H_{u,t}$。策略可以根据这些信息行动，或先将历史压缩为$\bm{h}_{u,t}$再计算$\pi(a_t\mid\bm{h}_{u,t})$。这种表示是否保留了未来预测所需的信息，是需要检验的建模条件。用循环网络读取历史，不会自动使压缩后的状态满足马尔可夫性。下文为了推导简洁写作$\pi(a_t\mid s_t)$；实际实现必须说明$s_t$由哪些可用信息替代，不能把模拟器的隐藏兴趣直接交给推荐策略。

“长期价值”还需要指定累计什么、累计多久。设一次轨迹在$T$时结束，折扣因子为$0\leq\gamma\leq1$，从$t$开始的回报为
\begin{equation}
 G_t=\sum_{\tau=t}^{T-1}\gamma^{\tau-t}r_\tau,
 \qquad
 J(\pi)=\mathbb{E}_{\pi}[G_0].
 \label{eq:rs-long-term-return}
\end{equation}
期望包含初始用户、策略抽样和环境转移的随机性。有限周期允许$\gamma=1$；无限周期常用$\gamma<1$并要求奖励有适当界限。固定两周的日历窗口和固定一百次请求是两个不同目标：后者可能只覆盖活跃用户，也可能因策略改变请求频率而对应不同的实际时长。若使用日历时间衰减，应根据相邻事件的真实时间间隔计算折扣，而不是把所有事件都当作等间隔。

奖励的单位决定了价值的含义。累计有效消费不等于满意度，累计点击不等于留存，平台净收益也不等于提供者收益。可以沿用第~\ref{chap:rs-constraints}~章的多目标权重，或把部分结果设为底线，但必须固定各分量口径。已经将一周内是否返回计作终点奖励时，不能又把同一个返回事件重复加入每天的奖励。观察窗口结束与用户真正离开也不同；前者是记录截断，后者在模型中可能是会影响未来回报的状态转移。

一个教学例子可以说明时间范围为何会改变选择。假设只剩两次请求，方案甲本次期望奖励为1，导致下一次期望奖励为0.2；方案乙本次为0.7，下一次为0.8。这里的第二次数值是假设的行动后果，后续策略相同，并非从当前点击率直接推出。当$\gamma=0.9$时，两方案价值分别为$1+0.9\times0.2=1.18$与$0.7+0.9\times0.8=1.42$。选择乙依赖这些转移估计可信；如果只观察到采用甲的日志，就不能仅凭这个算式证明乙更好。

从历史预测未来结果与估计改变行动后的结果也须分开。某类用户收到很多背包推荐后继续活跃，可能因为其原本需求强，也可能因为推荐确实帮助了他们。若日志未记录需求强度，条件预测会混合这两种解释。行动覆盖、可观测状态和额外干预决定了能够识别哪些后果，第~\ref{chap:rs-evaluation}~章将给出具体评价协议。连续决策的通用框架可参见\BookSectionRef{sec:reinforcement-learning}{强化学习}；本章保留推荐特有的行动、状态和证据要求。

图~\ref{fig:rs-long-term-feedback}把需要区分的三条路径放在同一个闭环中。系统选择曝光，改变了收集到的数据；曝光与消费也可能改变用户；推荐机会还可能改变提供者与后续目录。它们可以同时存在，但同一种日志变化未必能判定是哪条路径在起作用。

\begin{figure*}[htbp]
 \centering
 % 三条路径是待检验的机制；仅上路不要求用户偏好或供给发生变化。
\begin{tikzpicture}[
 x=1mm,y=1mm,font=\figurefont\small,
 block/.style={rectangle,draw=BookBlue,fill=BookBlue!5,
   line width=0.85pt,minimum width=34mm,minimum height=13mm,
   align=center,inner sep=2mm},
 endpoint/.style={block,minimum width=22mm,minimum height=17mm},
 flow/.style={-{Latex[length=2.2mm,width=1.5mm]},draw=BookInk,
   line width=1.1pt,rounded corners=1.2mm},
 wire/.style={draw=BookInk,line width=1.1pt,rounded corners=1.2mm},
 loop/.style={flow,draw=BookMuted,dashed,line width=0.85pt}
]
 \path[use as bounding box] (0,-45) rectangle (164,39);
 \node[endpoint] (action) at (12,0) {当前\\展示行动};
 \node[block] (data) at (55,26) {曝光与选择\\进入日志};
 \node[block] (train) at (106,26) {重新训练\\更新模型};
 \node[block,draw=BookAmber,fill=BookAmber!5] (interest) at (55,0)
   {接触与反馈\\改变兴趣或疲劳};
 \node[block,draw=BookAmber,fill=BookAmber!5] (user) at (106,0)
   {用户状态\\与可用历史};
 \node[block,draw=BookTeal,fill=BookTeal!5] (provider) at (55,-26)
   {机会与收益\\影响提供者};
 \node[block,draw=BookTeal,fill=BookTeal!5] (catalog) at (106,-26)
   {创建、修改、退出\\改变后续目录};
 \node[endpoint,draw=BookTeal,fill=BookTeal!5] (next) at (152,0)
   {下一轮\\可执行决策};

 \draw[wire] (action.east) -- (31,0);
 \draw[flow] (31,0) |- (data.west);
 \draw[flow] (31,0) -- (interest.west);
 \draw[flow] (31,0) |- (provider.west);
 \draw[flow] (data.east) -- (train.west);
 \draw[flow] (interest.east) -- (user.west);
 \draw[flow] (provider.east) -- (catalog.west);
 \draw[wire] (train.east) -- (132,26) -- (132,0);
 \draw[flow] (user.east) -- (next.west);
 \draw[wire] (catalog.east) -- (132,-26) -- (132,0);
 \draw[loop] (next.south) -- (152,-41) -- (12,-41) -- (action.south);

 \node[font=\figurefont\footnotesize,text=BookMuted] at (80,37)
   {数据选择循环：潜在偏好可以固定};
 \node[font=\figurefont\footnotesize,text=BookMuted,fill=white,
       inner sep=1mm] at (80,-41) {执行后继续下一轮交互};
\end{tikzpicture}

 \caption{连续推荐的三种反馈路径。上路只要求曝光影响数据与重训，中路另需用户兴趣或疲劳发生变化，下路另需提供者响应改变供给。箭头表示需要建模或检验的影响路径，不表示这些影响已经在真实系统中被识别；外侧回线表示下一轮决策。}
 \label{fig:rs-long-term-feedback}
\end{figure*}

\subsection{用户反馈与训练数据的共同演化}
\label{sec:rs-user-data-coevolution}

模型依据历史选择展示内容，新的交互又进入历史。即使用户潜在偏好保持不变，系统也会越来越多地看到自己曾经推荐的内容；若用户还会因接触而改变兴趣，就又增加了一条状态变化路径。要判断行为集中究竟来自数据选择还是偏好改变，必须先把这两个问题分开。

可以把一次重训写成$\theta^{(m+1)}=\operatorname{Train}(D^{(m)})$，其中$m$是训练轮次，$D^{(m)}$包括此前策略影响下采集的交互。决策时间$t$与训练轮次$m$不必一致，一轮训练可能跨越许多请求。只要$D^{(m)}$的组成依赖旧策略，新的$\theta^{(m+1)}$就会继承这种选择。下面两类诊断分别固定和放开潜在兴趣，用来判断不同反馈机制能够产生什么现象。

\Needspace{100mm}
\subsubsection{Algorithmic Confounding：重复训练的反馈诊断}
\label{sec:rs-algorithmic-confounding}

推荐记录不是对所有用户、所有物品的均匀测量。没有得到展示的物品通常没有后续选择记录，而重复训练可能把这种缺失造成的消费结构再次解释成偏好。\term{算法混杂}（algorithmic confounding）关注的正是这个问题：算法改变了训练数据的产生方式，却又将这些数据当作独立于算法的用户证据。Chaney等的研究用固定潜在效用的模拟环境隔离这一循环，比较只训练一次与持续收集交互并重新训练的推荐过程\scite{rs-e09}

在该环境中，用户有固定的潜在偏好，物品有固定属性，由此决定模拟中的真实效用。用户并不能在选择前完全知道效用；模型将效用区分为用户已知与尚未知晓的部分，用户的选择还受到推荐位置折扣影响。消费后才能获得更完整的反馈。每个用户对同一物品最多消费一次，新物品按外生规则进入目录。因此，已经消费过的对象、位置以及新增目录都会影响下一轮可观察的数据，不能把所有变化都归因于重训。

诊断的输入不是一份待去偏的评分矩阵，而是一整套环境设定：初始用户与物品、用户选择机制、目录到达过程、推荐模型和重训计划。给定这些条件，模拟器逐轮产生列表、选择与效用，再按照计划更新推荐模型。输出是多轮轨迹及其统计，而不是一种保证纠偏的排序器。比较时应让策略共享初始潜在偏好和外生到达序列，并报告不同重训计划何时能够访问新目录；否则，模型更新频率与目录可用性会一起变化。

例如，假设两个用户原本都同时喜欢登山包和相机，初始样本恰好使系统更常给两人展示登山包。新增交互矩阵会出现更多登山包选择，重新拟合的表示可能进一步靠近登山用品。这个过程不要求两人的真实喜好发生任何变化。为了判断它是否带来损失，不能只看交互是否更一致，还要利用模拟器已知的真实效用，检查推荐是否错过了更合适的对象。

行为相似性可以按两人消费集合$B_u$与$B_v$的Jaccard系数衡量：
\[
 \operatorname{sim}(u,v)
 =\frac{|B_u\cap B_v|}{|B_u\cup B_v|}.
\]
假设两人都只消费了登山包和水壶，这个值为1；它说明消费集合相同，并未测量潜在偏好距离。原研究还按学到的用户表示选取近邻，再与相同用户配对在理想推荐条件下的交互作比较。保留相同配对很重要，否则“近邻更相似”可能只是比较对象已被换过。

固定偏好的设定使一个解释得到排除：模拟中的行为同质化不能归因于潜在偏好本身被修改。剩余机制仍可能包括曝光选择、位置影响和已消费对象退出可选范围，必须结合对照定位。该研究在其环境下观察到的效用与同质化变化，不是所有推荐系统必然退化的定理。真实平台没有可直接读取的潜在效用表；要复现同样的解释力度，需要随机曝光、可靠的偏好测量或其他能够区分替代解释的证据。

\Needspace{100mm}
\subsubsection{Degenerate Feedback Loops：兴趣变化的动态诊断}
\label{sec:rs-degenerate-feedback}

固定偏好能够隔离数据选择，却不能描述用户在接触内容后发生的变化。例如，连续阅读某一主题可能增强兴趣，也可能因重复而疲劳。Jiang等的\term{退化反馈循环}研究将兴趣本身作为动态状态，分析推荐、响应和兴趣更新之间的循环。在这里，退化（degeneracy）有一个特定数学含义：兴趣状态持续远离初始状态，不是直接以低满意度或低社会福利定义\scite{rs-e10}

在固定有限目录$I$上，记$\bm{\xi}_t=(\xi_{i,t})_{i\in I}$为兴趣向量。若几乎处处有
\[
 \limsup_{t\to\infty}
 \lVert\bm{\xi}_t-\bm{\xi}_0\rVert_2=\infty,
\]
可称为弱退化；将$\limsup$换为$\lim$则是强退化。前者允许反复回到较近的状态，只要求存在越来越远的时点；后者要求最终持续远离。定义依赖兴趣数值的尺度与允许范围。如果兴趣被模型强制限制在紧集内，这种无界偏离的定义便不可能满足，但仍可能出现其他意义上的集中或体验损失。

为看清更新机制，作一个简化教学假设：本次只曝光物品$i$，点击$y_t\in\{0,1\}$服从$\sigma(\xi_{i,t})$，其中$\sigma$为逻辑函数；兴趣按正步长$\delta$更新为
\begin{equation}
 \xi_{i,t+1}=\xi_{i,t}+\delta(2y_t-1),
 \qquad
 \xi_{j,t+1}=\xi_{j,t}\quad(j\ne i).
 \label{eq:rs-interest-reinforcement}
\end{equation}
这表示点击后兴趣增强、未点击后减弱。它是便于演示的正强化模型；原研究的模拟还允许不同对象具有不同方向的更新参数，不能把此处结论套给所有参数。

若$\xi_{i,0}=0$、$\delta=0.2$，前四次曝光的响应恰好为$1,1,1,0$，状态依次为$0.2,0.4,0.6,0.4$，对应下一次点击概率约为$0.550,0.599,0.646,0.599$。一次未点击会回退，但是否长期发散取决于完整随机过程。条件平均增量为
\[
 \mathbb{E}[\xi_{i,t+1}-\xi_{i,t}\mid\xi_{i,t}]
 =\delta\bigl(2\sigma(\xi_{i,t})-1\bigr).
\]
正兴趣下平均增量为正，负兴趣下为负，这说明局部具有远离零点的倾向。它本身还不是几乎处处发散的证明；无限时间结论另需兴趣更新、噪声和曝光频率等条件。

由此也能理解，扩大探索范围与稳定兴趣幅度是两回事。在有限目录上，即使每次均匀随机推荐，各对象仍可能获得无限多次曝光；若其更新规则具有足够强的强化机制，随机探索并未切断变化过程。目录持续增长则改变了曝光分配，单个对象未必继续获得同样多的机会，但是否避免退化仍与增长速度和策略有关。

这类诊断应同时报告曝光覆盖和兴趣状态轨迹。只观察到覆盖更多对象，不能认定兴趣动态稳定；只观察到兴趣范数增大，也不能不经测量就认定用户满意度下降。原研究提供的是特定动态条件下的分析和模拟比较，说明哪些反馈规则值得警惕，而不是一种通用的消除信息茧房算法。若实际任务主要表现为重复后的疲劳，应改用能表达恢复、饱和或厌倦的转移模型，并重新检查结论。

\subsection{供给者响应与目录演化}
\label{sec:rs-provider-dynamics}

前两条路径都可以在供给外生的条件下讨论。现实中，推荐机会还会影响提供者是否继续投入、生产什么以及面向谁生产。因而时点$t+1$的目录$I_{t+1}$可能是当前推荐的后果，而不只是第~\ref{sec:rs-candidate-maintenance}~节所处理的一次外部数据更新。

设提供者集合为$C$，$c(i)$表示物品$i$所属提供者。记$\bm{z}_{c,t}$为其决策前状态，可包含近期有效消费、收入、生产成本、库存或继续生产的意愿。推荐产生机会与反馈$e_{c,t}$，提供者据此选择生产或修改动作$b_{c,t}$。一种抽象转移关系为
\[
 \bm{z}_{c,t+1}\sim
 P_c(\,\cdot\mid\bm{z}_{c,t},e_{c,t},b_{c,t}),
 \qquad
 I_{t+1}=\operatorname{Update}(I_t,\{b_{c,t}\}_{c\in C}).
\]
这里的$\operatorname{Update}$可以包含新建、修改和退出，不是一项固定算法。提供者反应可能按周发生，而推荐按秒发生；将许多次曝光聚合到同一提供者时间步时，须保留聚合窗口，不能在每次请求后虚构一次完整的创作响应。

相同曝光也不意味着相同激励。假设甲、乙各获得100次曝光，甲有10次有效消费，每次净收益1元，本轮固定生产成本8元；乙有4次有效消费，每次净收益3元，固定成本2元。按这个教学口径，甲净得2元，乙净得10元。如果下一轮投入门槛是5元，两者可能作出不同的继续生产决定。差异来自消费、收益和成本，而不是曝光数量本身；门槛规则也是例子假设，不能据此预测真实商家。

提供者效用、存续、内容质量和内容差异因此需要分开记录。提高总效用可能主要帮助已有大提供者，也可能改变小提供者的退出概率；增加目录数量又可能只是同类物品重复增加。静态曝光公平约束给出了机会分配标准，却没有自动指定机会怎样转化为未来供给。

供给还会连接不同用户。一位提供者根据一组用户的推荐机会调整内容，调整后的目录可能被另一组用户看到。此时逐用户独立的转移近似会遗漏跨用户影响。后面的EcoAgent将提供者累计效用放进策略目标，Performative Recommendation则预判评分规则引起的内容修改；两者都需要明确响应模型，其模拟结论不能直接代替真实市场中的进入、退出和竞争证据。

\section{连续决策学习}
\label{sec:rs-sequential-learning}

问题定义以后，可以学习两类对象：价值函数回答在指定后续策略下会得到多少累计结果，策略回答在当前信息和约束下怎样选择。二者能够交替更新，也能够联合学习；有的方法直接从轨迹改策略，不先学习显式环境模型。共同要求是价值、行动和执行范围相容，不能用一个列表的估值去评价另一个经过任意去重或重排的列表。

\subsection{长期价值估计}
\label{sec:rs-long-term-estimation}

给定策略$\pi$，状态价值$V^\pi(s)$表示从状态$s$开始继续执行$\pi$的期望回报；行动价值$Q^\pi(s,a)$则先固定当前行动$a$，之后再执行$\pi$。在状态充分、时间范围已进入状态或终止规则的条件下，
\begin{equation}
 \begin{aligned}
 Q^\pi(s,a)
   &=\mathbb{E}[r+\gamma V^\pi(s')\mid s,a],\\
 V^\pi(s)
   &=\mathbb{E}_{a\sim\pi(\cdot\mid s)}[Q^\pi(s,a)].
 \end{aligned}
 \label{eq:rs-policy-value}
\end{equation}
把未来价值接在当前奖励之后，就能让一个暂时无收益但改善后续条件的行动获得信用。该递推的通用推导见\BookSectionRef{sec:decision-dynamic-programming}{序贯策略与动态规划}。

直接价值学习可以用完整回报$G_t$作为目标，也可以用$r_t+\gamma\widehat V(s_{t+1})$作时序差分目标。前者等待轨迹成熟，后者使用当前估计补上尚未观察到的后果。学习环境模型则先估计奖励与转移，再在模型中展开后续行动。两条路径都能输出价值，但错误来源不同：直接估值受日志覆盖与目标估计影响，模型展开还会累积转移误差。训练损失较小，只说明在训练分布及其目标下拟合较好。

终止与延迟也直接影响训练样本。真正的终止状态可令后续价值为零；日志文件结束却不表示用户此后没有价值，此时应明确截断处理或补取后续记录。转化延迟跨越多次请求时，还须说明奖励分配给哪个事件。若只保留已消费记录而丢掉未消费后的离开，得到的训练分布会系统性缺少一种行动后果。

推荐中的额外困难来自组合行动。候选数为$n$、有序列表长度为$k$时，不重复列表有$n!/(n-k)!$种。为每种列表分别学习价值不仅需要大量数据，也难以复用物品之间的结构。下面的SlateQ利用用户每次实际消费的对象分解价值，减少价值学习规模；列表内部竞争仍由选择概率保留。

\Needspace{110mm}
\subsubsection{SlateQ：列表价值分解}
\label{sec:rs-slateq}

列表包含多个对象，但一次交互可能只消费其中一个。\term{SlateQ}据此引入\term{条件物品价值}：把实际消费对象固定以后，这次消费及后续策略带来的期望回报。这里的“物品价值”不是该物品独立展示时的价值，而是在明确条件下能够跨列表复用的价值分量\scite{rs-r18}

记$c_t$为实际消费对象，$c_t=0$表示没有消费；$p(i\mid s,L)$表示在状态$s$下看到列表$L$后选择$i$的概率。分解依赖两个不同条件：
\begin{itemize}
 \item 每一步至多消费一个对象，也允许不消费，因此$c_t\in L\cup\{0\}$，这些互斥事件的概率和为1。
 \item 给定状态$s$与实际消费对象$i$后，期望奖励和下一状态分布不再依赖列表的其他内容。
\end{itemize}
第一项只限制消费数量，不蕴含第二项。例如，用户只点击一个背包，仍可能因旁边反复出现的相机广告而疲劳。第二项也不要求$p(i\mid s,L)$与其他对象无关；恰恰相反，物品能否被选中仍取决于完整列表及位置。

在第二项条件下，可以定义与$L$无关的条件价值
\begin{equation}
 \overline Q^\pi(s,i)
 =\mathbb{E}[r+\gamma V^\pi(s')\mid s,c=i].
 \label{eq:rs-slateq-item-value}
\end{equation}
这里的条件分布由假设中可跨列表复用的奖励与转移确定。对实际消费事件使用全期望公式，就得到
\begin{equation}
 Q^\pi(s,L)=
 \sum_{i\in L\cup\{0\}}
 p(i\mid s,L)\overline Q^\pi(s,i).
 \label{eq:rs-slateq-decomposition}
\end{equation}
分解的关键不是把列表看作一组互不相关的物品，而是把“会选择谁”和“选择之后会怎样”分开。前者仍是列表问题，后者才在条件成立时降为物品问题。

零消费项不能省略。用户本次没有消费，仍可能继续浏览、结束会话，或在后续请求得到帮助，因此$\overline Q^\pi(s,0)$一般不为零。只有奖励和后续价值都符合相应条件时才可将其置零。若不消费时是否离开取决于未点击列表的具体内容，零消费这一分支已经违反第二项条件。

图~\ref{fig:rs-long-term-slateq}给出一个教学计算。假设两个对象与零消费的选择概率为$0.5,0.3,0.2$，条件价值为$4,2,0$，则列表价值为
\[
 0.5\times4+0.3\times2+0.2\times0=2.6.
\]
如果仅把零消费价值改为1，其他量不变，列表价值便是2.8。这个变化不是增加一次点击，而是对未消费后续路径作出了不同评价。

\begin{figure}[htbp]
 \centering
 % 教学数值；条件价值分支互斥，零消费分支不能在一般情形下删除。
\begin{tikzpicture}[
 x=1mm,y=1mm,font=\figurefont\small,
 block/.style={rectangle,draw=BookBlue,fill=BookBlue!5,
   line width=0.85pt,minimum width=34mm,minimum height=13mm,
   align=center,inner sep=1.5mm},
 flow/.style={-{Latex[length=2.2mm,width=1.5mm]},draw=BookInk,
   line width=1.1pt,rounded corners=1.2mm},
 wire/.style={draw=BookInk,line width=1.1pt,rounded corners=1.2mm},
 prob/.style={font=\figurefont\footnotesize,text=BookInk,inner sep=1mm}
]
 \path[use as bounding box] (0,-7) rectangle (112,61);
 \node[block,minimum width=22mm] (slate) at (12,27) {状态$s$\\列表$(1,2)$};
 \node[block] (one) at (58,50) {消费物品1\\$\overline Q^\pi(s,1)=4$};
 \node[block] (two) at (58,27) {消费物品2\\$\overline Q^\pi(s,2)=2$};
 \node[block,draw=BookMuted,fill=BookPaper] (none) at (58,4)
   {不消费\\$\overline Q^\pi(s,0)=0$};
 \node[block,minimum width=24mm,draw=BookTeal,fill=BookTeal!5]
   (value) at (99,27) {概率加权\\$Q^\pi=2.6$};

 \draw[wire] (slate.east) -- (30,27);
 \draw[flow] (30,27) |- (one.west);
 \draw[flow] (30,27) -- (two.west);
 \draw[flow] (30,27) |- (none.west);
 \node[prob,anchor=south] at (35,50) {$0.5$};
 \node[prob,anchor=south] at (35,27) {$0.3$};
 \node[prob,anchor=south] at (35,4) {$0.2$};
 \draw[wire] (one.east) -- (81,50) -- (81,27);
 \draw[wire] (none.east) -- (81,4) -- (81,27);
 \draw[flow] (two.east) -- (value.west);
\end{tikzpicture}

 \caption{SlateQ的条件分解。列表影响各消费分支的概率；给定消费对象后，条件价值在假设成立时不再依赖其他展示对象。图中的概率与价值是教学假设，零消费价值为零仅属于本例。}
 \label{fig:rs-long-term-slateq}
\end{figure}

学习时，一条记录至少包含$(s_t,L_t,c_t,r_t,s_{t+1})$以及终止信息。选择模型可由包括零消费在内的选择事件拟合；条件价值则用实际发生的分支更新。若下一列表$L_{t+1}$按当前策略$\pi$执行，一种采样时序差分目标为
\[
 y_t=r_t+\gamma
 \sum_{j\in L_{t+1}\cup\{0\}}
 \widehat p(j\mid s_{t+1},L_{t+1})
 \widehat{\overline Q}^{\,\pi}(s_{t+1},j).
\]
在表格表示中，只更新当前访问的$(s_t,c_t)$：
\[
 \widehat{\overline Q}^{\,\pi}(s_t,c_t)
 \leftarrow
 \widehat{\overline Q}^{\,\pi}(s_t,c_t)
 +\eta\bigl(y_t-\widehat{\overline Q}^{\,\pi}(s_t,c_t)\bigr),
\]
其中$\eta$为学习率。使用参数网络时，可以最小化当前分支预测与停止梯度的$y_t$之间的平方误差。假设当前预测为4，实际奖励为1，下一列表估值为3，$\gamma=0.9$、$\eta=0.1$，则$y_t=3.7$，更新后为3.97。这是在修正一次条件价值估计，不是直接把整个列表的分数改成3.97。

如果目的是最优价值学习，则下一状态的续接项应改为合法列表上的最大值，
\[
 y_t=r_t+\gamma\max_{L\in\mathcal{F}_{t+1}}
 \sum_{j\in L\cup\{0\}}
 \widehat p(j\mid s_{t+1},L)
 \widehat{\overline Q}(s_{t+1},j).
\]
当前策略评价与这种Q学习目标的后续策略不同，不能混用后仍把结果解释为同一个$Q^\pi$。最大化怎样求解也会影响训练目标，近似列表求解的误差会进入价值学习。

分解减少了价值表示的规模，并未消除状态缺失和选择模型错误。补充决策前累计频次，有时能够修复遗漏疲劳历史的问题；但若当前未消费对象本身改变未来兴趣，仅补历史无法使其当前影响消失。此时需要扩大消费事件或行动建模范围，乃至直接建模完整列表后果。SlateQ原研究中的参与目标证据，也不能越过奖励定义被解读为对长期满意度或福利的普遍保证。

\subsection{连续策略优化}
\label{sec:rs-sequential-policy}

价值估计确定了评价对象，策略优化还须决定哪些行动值得执行。可以用估计价值在候选列表中搜索，也可以直接提高带来较高回报的行动概率；频次控制、主动询问和试探冷启动对象，则扩展了行动空间。无论采用哪条路径，都要沿用相同的资源账本和展示规则。

探索的长期价值来自信息会影响后续选择。作一个两次机会的教学假设：已知物品每次奖励都是0.6；新物品以相同概率属于奖励恒为1或恒为0的两种类型，第一次尝试即可无误识别类型。始终展示已知物品的总奖励为1.2。先尝试新物品，第二次在好类型下继续展示它，在差类型下改回已知物品，期望总奖励为$0.5+0.5\times1+0.5\times0.6=1.3$。这种优势依赖可学习的类型和剩余机会；只有一次机会时，试探反而将期望奖励从0.6降到0.5。

信息价值的通用计算见\BookSectionRef{sec:decision-value-of-information}{信息价值与探索}，推荐中的探索曝光与体验代价见第~\ref{sec:rs-exploration-exposure}~节。实际问题没有例子中的无误识别，试探还可能消耗预算或引起疲劳，需要将这些后果计入同一个累计目标。执行之后，应记录实际行动、决策时的可见信息及相应抽样概率；仅保存优化器原拟议的列表，无法重建后续过滤、去重或重排产生的分布。

\subsubsection{基于长期价值的列表改进}
\label{sec:rs-long-term-slate-improvement}

在SlateQ条件成立时，给定$\overline Q^\pi$、选择模型和可行列表，可以寻找相对于当前策略更好的本次行动：
\begin{equation}
 L^+(s)\in\argmax_{L\in\mathcal{F}(s)}
 \sum_{i\in L\cup\{0\}}
 p(i\mid s,L)\overline Q^\pi(s,i).
 \label{eq:rs-long-term-slate-improvement}
\end{equation}
这里先固定后续策略$\pi$。在精确的马尔可夫模型、相同折扣目标和完整可行范围下，逐状态选择不低于原策略价值的行动，可以使用策略改进结论；若状态、价值或最大化都只是近似，实际效果仍须另验。

表~\ref{tab:rs-long-term-slate-example}列出一个三候选、二位置的教学例子。物品1、2、3分别为登山包、水壶和相机，条件价值设为$4,2,3.5$，零消费价值为0。本例规定组合内按编号顺序展示，因而恰好只有三种合法列表；允许任意位置交换时，还要分别评估其他顺序。

\begin{table}[htbp]
 \centering
 \small
 \begin{tabularx}{\linewidth}{@{}lXrr@{}}
  \toprule
  列表 & 对应物品、零消费概率 & 消费概率 & 长期价值 \\
  \midrule
  $(1,2)$ & $(0.50,0.30,0.20)$ & 0.80 & 2.600 \\
  $(1,3)$ & $(0.20,0.60,0.20)$ & 0.80 & 2.900 \\
  $(2,3)$ & $(0.25,0.65,0.10)$ & 0.90 & 2.775 \\
  \bottomrule
 \end{tabularx}
 \caption{相同条件物品价值下的列表枚举。每行概率按该行列表顺序给出，最后一项为零消费。消费即点击仅为本例的即时比较口径；所有数值均为教学假设。}
 \label{tab:rs-long-term-slate-example}
\end{table}

若只最大化发生消费的概率，会选$(2,3)$。长期价值则分别为$0.5\times4+0.3\times2=2.6$、$0.2\times4+0.6\times3.5=2.9$和$0.25\times2+0.65\times3.5=2.775$，因而选择$(1,3)$。这里并不是把点击分数换成物品价值后独立排序：同一物品1在两个列表中的选择概率从0.5变为0.2，竞争关系参与了结果。

独立取条件价值最大的$k$个对象也不普遍正确。另设一个教学反例，三者价值为$10,9,8$，零消费价值为0；列表$(1,2)$的概率为$(0.05,0.05,0.90)$，列表$(1,3)$的概率为$(0.40,0.40,0.20)$。取最大的两个物品只得$0.95$，后一列表得到$7.2$。这些概率满足基本归一化条件，已足以反驳不附加选择模型条件的独立排序保证。

候选规模较小时，枚举能够直接核对合法性与最优值。规模变大后，可以使用贪心、局部搜索或受限候选搜索，但必须说明搜索范围与停止条件，不能把“算出了一个列表”写成全局最优。若选择模型具有特定形式，例如$p(i\mid s,L)=w_i/(w_0+\sum_{j\in L}w_j)$，则目标变为
\[
 \frac{w_0\overline Q^\pi(s,0)
       +\sum_{i\in L}w_i\overline Q^\pi(s,i)}
      {w_0+\sum_{i\in L}w_i}.
\]
分子、分母都随列表变化，可以研究对应的分式优化；这仍不等于直接按$\overline Q^\pi$取前$k$个。若选择权重随位置变化，或者还要求商家配额，求解器也必须保留这些变量和约束。

执行改进列表后，新轨迹会改变估值所依据的分布。可在一段时间内固定策略再评价，或明确使用持续改进的更新目标。日志中某列表估值更高，只是候选改进的依据；覆盖不足、选择模型偏差与上线后的状态变化都可能改变排序，不能省略后续效果检验。

\subsubsection{Top-K REINFORCE：采样校正的列表策略学习}
\label{sec:rs-topk-reinforce}

当目录很大时，每次对组合列表求最大值可能过于昂贵。另一条路径是直接学习物品抽样分布，用累计回报提高有利选择的概率。\term{Top-K REINFORCE}把这种策略梯度方法用于大规模推荐候选生成，并处理日志策略与目标策略不同以及一次输出多个候选的问题\scite{rs-r91}

模型根据用户信息或循环网络编码的交互历史产生状态表示$\bm{h}_{u,t}$，再与物品向量计算分数。为强调其策略含义，记单次抽样分布为
\[
 p_\theta(i\mid\bm{h})
 =\frac{\exp(\langle\bm{h},\bm{v}_i\rangle)}
        {\sum_{j\in I_t}\exp(\langle\bm{h},\bm{v}_j\rangle)}.
\]
参数$\theta$包含状态网络和物品参数。物品很多时，原实现使用采样softmax近似训练计算；使用近似归一化会引入相应的估计问题，不能把未归一化检索分数直接当作已知概率。

在同策略轨迹上，REINFORCE用“回报乘对数概率梯度”调整选择。以有限轨迹目标为例，梯度具有
\[
 \mathbb{E}_{\pi_\theta}
 \left[\sum_t\gamma^t G_t
       \nabla_\theta\log\pi_\theta(a_t\mid\bm{h}_{u,t})\right]
\]
的形式；$\gamma=1$时外侧折扣消失。高回报行动得到更大的概率提升，低回报或低于基线的行动则相反。状态表示、动作概率与累计反馈因此通过同一条轨迹联系起来，而非逐条把点击作为独立正例。

日志来自行为策略$\mu$时，轨迹的状态和行动分布都可能不同。在环境转移机制相同且有覆盖的条件下，完整轨迹的重要性权重包含各时点行动比值的乘积。乘积容易有很大方差，原方法将修正近似为当前步的目标概率与行为概率之比。这个简化没有完整校正早先行动造成的状态分布变化，因而是有偏近似，不是任意离策略日志上的无偏策略梯度。

行为概率也通常需要估计。原方法另设行为头$\widehat\mu(i\mid\bm{h})$，共用状态表示，却在该分支阻断流向状态表示的梯度，以免行为拟合目标直接改变策略状态网络。它使用全部相应日志记录拟合行为分布，而不是只看非零奖励的记录。只从成功记录估计“当时有多可能选到这个物品”，会把奖励选择又混进分母。即使采用行为头，概率估计误差、小概率比值以及日志缺失仍会影响训练。

这里的Top-K也有准确的抽样语义。方法从$p_\theta$独立有放回抽样$K$次，再删除重复物品；最终集合大小可能小于$K$。某物品至少出现一次的包含概率为
\begin{equation}
 \alpha_\theta(i\mid\bm{h})
 =1-\bigl(1-p_\theta(i\mid\bm{h})\bigr)^K,
 \qquad
 \frac{\partial\alpha_\theta}{\partial p_\theta}
 =K(1-p_\theta)^{K-1}.
 \label{eq:rs-topk-inclusion}
\end{equation}
这不是按分数确定性取前$K$名，也不是恰好抽出$K$个不同对象。各物品包含概率之和是期望去重后集合大小，通常不为1，不能把$\alpha_\theta$再当作普通物品分布解释。

在可加物品收益的近似目标下，简记当前日志物品$i_t$的概率为$p_t=p_\theta(i_t\mid\bm{h}_{u,t})$、$\widehat\mu_t=\widehat\mu(i_t\mid\bm{h}_{u,t})$和$\alpha_t=1-(1-p_t)^K$。对应的一步修正可以写成
\begin{equation}
 \begin{aligned}
 \widehat{\bm{g}}_t
 &=G_t\,\frac{\alpha_t}{\widehat\mu_t}
   \nabla_\theta\log\alpha_t\\
 &=G_t\,\frac{p_t}{\widehat\mu_t}
   K(1-p_t)^{K-1}\nabla_\theta\log p_t.
 \end{aligned}
 \label{eq:rs-topk-policy-gradient}
\end{equation}
此处省略了轨迹外侧折扣或抽样加权，只展示当前记录的修正；行为概率估计在策略更新中视为固定。等式由链式法则得到。不能先乘包含概率比，又额外重复乘一次导数因子，那会改变原梯度。

假设一段教学轨迹的两步奖励为0和2，$\gamma=0.9$，则首步$G_t=1.8$。若该步$p_\theta=0.2$、$\widehat\mu=0.1$、$K=2$，包含概率是0.36，导数因子是1.6。相对于$\nabla_\theta\log p_\theta$，式~\eqref{eq:rs-topk-policy-gradient}的系数为$1.8\times2\times1.6=5.76$。若单次概率为0.8，同样$K=2$时导数因子只有0.4：已经很容易进入集合的物品，继续增加单次概率对包含率的边际作用较小。

可加收益是此处的重要边界。若集合收益可写成$\sum_{i\in L}g(s,i)$，其期望才可写为$\sum_i\alpha_\theta(i\mid s)g(s,i)$；去重后每个物品只贡献一次。SlateQ则将至多一个实际消费的互斥分支加权，两者的概率和与收益语义都不同。即时收益可加也不能推出任意长期状态转移可加，尤其不能自动处理位置竞争、重复疲劳或供给联动。

原论文的在线场景中，该模型是多路候选生成的一路，下游还有排序；累计回报采用小时级窗口，文中使用的范围为4至10小时。因此，其证据不能直接升级为长期留存改善，也不能把候选采样概率当作最终页面曝光概率。将它用于更长周期时，需要重新定义状态、奖励、行为概率和覆盖范围，并检验下游选择怎样改变策略实际产生的后果。

\subsubsection{EcoAgent：用户与提供者效用的连续优化}
\label{sec:rs-ecoagent}

仅根据用户回报学习，可能忽略本次机会对未来内容供给的影响。\term{EcoAgent}把用户与提供者都作为有历史、有状态的主体，学习它们的累计效用，再据此优化推荐策略。其提供者部分关注的不是谁目前总价值最高，而是当前一次推荐相对于不推荐，会为被选提供者增加多少后续效用\scite{rs-r119}

记$\bm{h}_{u,t}$为用户历史编码，$\bm{z}_{c,t}$为提供者历史编码。用户反馈可按连续交互输入循环网络；提供者反馈在较慢时间尺度上聚合，例如一段时间内获得的机会、有效消费与用户反馈。两类价值网络分别拟合用户回报和提供者回报，原方法采用Huber损失，以降低较大残差对拟合的过度影响。策略网络则同时读取用户状态、候选物品及其提供者状态，输出推荐分布。

为解释增量从何而来，先考虑每次推荐一个物品$i$的简化行动。把用户条件累计效用记为$Q_u^\pi(s,i)$，提供者$c$在该行动及后续策略下的累计效用记为$Q_c^\pi(s,i)$。在固定量纲和权重$\lambda\in[0,1]$后，总行动价值可写为
\[
 Q_{\mathrm{tot}}^\pi(s,i)
 =(1-\lambda)Q_u^\pi(s,i)
   +\lambda\sum_{c\in C}Q_c^\pi(s,i).
\]
这一式子仍要对所有提供者求和。若每次推荐都有复杂竞争影响，不能只计算被选提供者便声称优化了该总目标。

EcoAgent所用的简化要求：在当前状态与后续策略固定时，未被推荐者的累计后果不因“究竟推荐了另一个谁”而改变。设$Q_{c,0}^\pi(s)$表示提供者$c$本次未获推荐的条件累计效用，并定义被选者增量
\[
 \Delta_c^\pi(s,i)
 =Q_c^\pi(s,i)-Q_{c,0}^\pi(s),
 \qquad c=c(i).
\]
如果未选者满足上述不干扰条件，就有
\begin{equation}
 \sum_{c\in C}Q_c^\pi(s,i)
 =\underbrace{\sum_{c\in C}Q_{c,0}^\pi(s)}_{B^\pi(s)}
  +\Delta_{c(i)}^\pi(s,i).
 \label{eq:rs-provider-increment}
\end{equation}
共同基线$B^\pi(s)$在条件于$s$的策略梯度中消失，因为
$\sum_i\pi_\theta(i\mid s)\nabla_\theta\log\pi_\theta(i\mid s)=0$。
于是当前更新可以使用
\[
 (1-\lambda)Q_u^\pi(s,i)
 +\lambda\Delta_{c(i)}^\pi(s,i)
\]
作为与行动相关的价值部分。这里消失的是同一状态、同一后续策略下的基线，不是对所有策略都相同的全局常数；策略改变后，状态分布与$B^\pi$也会改变。

这个条件涉及累计后果，强于“本次没有直接改动其他提供者的收入”。即使本次只给甲一次机会，甲之后增加供给也可能占用乙未来的曝光；后续策略还会根据甲的新状态重新分配流量。这样的跨提供者影响可能使$Q_{c,0}^\pi(s)$依赖当前被选者身份，式~\eqref{eq:rs-provider-increment}便不再成立。应用时应说明环境与未来策略为何允许该简化，或显式估计其他提供者的变化。

增量本身需要反事实估计。原方法通过从提供者历史中移除当前一次推荐，构造“本次未获推荐”的输入，再用同一个效用模型分别评价事实历史和修改后的历史，二者之差用于近似增量。真实日志只记录了实际发生的轨迹，并没有同时观察这个反事实。模型是否见过类似状态、移除事件后后续输入是否仍合理，都影响增量可靠性；由网络输出两个数相减，不会自动得到可识别的因果效应。

一个教学例子可以解释累计水平与边际机会为何不同。假设用$U(n)=\log(1+n)$描述某提供者已获$n$次有效机会时的价值，并暂不考虑未来竞争。甲已有9次，乙已有1次，再增加一次的增量分别为
\[
 \Delta_{\text{甲}}=\log(11/10)\approx0.0953,
 \qquad
 \Delta_{\text{乙}}=\log(3/2)\approx0.4055.
\]
甲的当前价值$\log10$高于乙的$\log2$，但乙的下一次机会更有边际作用。若用户侧价值分别为0.8和0.6，$\lambda=0.5$，组合值约为0.4477和0.5027，此时选择乙；若$\lambda=0.2$，组合值约为0.6591和0.5611，则选择甲。这个例子使用假设的饱和曲线，不是原论文的效用测量，也没有保证少机会者总应优先。

训练由与当前策略交互的轨迹提供依据。系统收集用户与提供者反馈，形成两类成熟回报并拟合价值网络，再计算提供者事实与反事实价值，用REINFORCE更新选择概率。策略梯度使用价值估计时按相应估计量处理，不把反事实网络的所有输入改动都当作真实环境的可微导数。更新后的策略继续产生新轨迹，因此价值拟合与策略改进可以多轮交替。

上线时的直接输出仍是可以执行的推荐分布或物品选择。选择乙会更新乙的机会和用户反馈，是否进一步增加供给由环境响应决定，不能在策略得分里预设必然发生。原研究主要在模拟生态中检验这些联系，而且指出提供者总效用提高时，存续人数仍可能减少。总效用、供给差异与存续必须分别评价；“生态”这个目标名称不能代替对具体后果的核验。

\subsubsection{Performative Recommendation：创作响应的策略建模}
\label{sec:rs-performative-recommendation}

机会影响供给数量之外，评分规则还会影响内容朝哪个方向修改。若相机和背包内容都为了迎合相同受众而趋向同一种表达，固定目录上的多样性重排只改变当前分配，没有描述创作者下一步的选择。\term{Performative Recommendation}研究会反过来改变被评分对象的推荐：评分模型公开或被适应之后，提供者根据激励修改内容，模型应预判这种响应\scite{rs-r120}

原模型以单位范数向量表示用户评分方向$\bm{v}_u$和物品当前特征$\bm{x}_i$。每个物品$i$有一组可以触达的候选用户$U_i$，评分采用内积。将这些用户评分向量的平均方向归一化为$\overline{\bm{v}}_i$后，创作者选择新的单位向量$\bm{x}'_i$，提高受众评分，同时支付偏离原内容的二次成本：
\begin{equation}
 \bm{x}'_i\in
 \argmax_{\lVert\bm{x}'\rVert_2=1}
 \left\{\langle\overline{\bm{v}}_i,\bm{x}'\rangle
       -\alpha\lVert\bm{x}'-\bm{x}_i\rVert_2^2\right\}.
 \label{eq:rs-performative-response-objective}
\end{equation}
这里$\alpha\geq0$是修改成本系数，不是反馈学习率。可触达用户集合也不是任意可改变量，它限制了评分规则能为创作者提供哪些激励。

由于原向量和新向量都是单位向量，二次项等于
$2\alpha-2\alpha\langle\bm{x}',\bm{x}_i\rangle$。
去掉与$\bm{x}'$无关的常数，问题成为最大化
$\langle\overline{\bm{v}}_i+2\alpha\bm{x}_i,\bm{x}'\rangle$，
因此非零情形下的最佳响应为
\begin{equation}
 \bm{x}'_i=
 \frac{\overline{\bm{v}}_i+2\alpha\bm{x}_i}
      {\lVert\overline{\bm{v}}_i+2\alpha\bm{x}_i\rVert_2}.
 \label{eq:rs-performative-response}
\end{equation}
这个表达式说明，修改后方向在受众评分方向与原内容方向之间取舍。若分母为零，目标在单位球上不再给出唯一选择，需要另定平局响应，不能直接做归一化；候选用户平均向量为零时，前面的平均方向也需单独约定。

假设两个物品初始特征为$(1,0)$与$(0,1)$，都面向平均评分方向$(1,0)$的相同受众。在$\alpha=0$的教学极端情形下，两者的最佳响应都变成$(1,0)$，初始余弦距离为1，修改后为0。若$\alpha=1$，第一个保持$(1,0)$，第二个变为$(1,2)/\sqrt5$，两者余弦距离约为$1-1/\sqrt5=0.5528$。有限修改成本保留了一部分差异，但不能由这一例子断言所有成本设置都能维持多样性。

训练评分模型时，可以把式~\eqref{eq:rs-performative-response}接入目标，让参数更新看到预测的内容变化。需要特别分清原研究中两个目标的时间口径：排序质量使用修改前数据上的相关性标签与NDCG，多样性使用预测修改后的内容。用矩阵$\bm{X}$汇集当前物品特征，$\bm{Y}$保存当前用户与物品的相关性标签，$\bm{X}'(\theta)$表示响应映射给出的修改后特征，其代理目标可概括为
\begin{equation}
 \begin{aligned}
 \max_\theta\ \bigl\{&
 \widetilde{\operatorname{NDCG}}(\theta;\bm{X},\bm{Y})\\
 &+\lambda\,\widetilde{\operatorname{Div}}
       (\theta;\bm{X}'(\theta))\bigr\}.
 \end{aligned}
 \label{eq:rs-performative-learning}
\end{equation}
其中$\lambda$为取舍权重，波浪号表示可微近似。原方法用软排序与软Top-$k$近似离散排名和集合选择，使梯度能够经过评分、响应映射和多样性计算。不能把第一项改述成“已经知道修改后内容对用户的真实质量”：未来相关性标签尚未被观察，这正是连续部署后需要重新收集的信息。

多样性在此来自推荐规则改变了创作激励，而不只是给当前列表添加更多类别。假如两个创作者只能面对完全相同的受众，且修改成本很低，评分规则可实现的差异化激励就可能受限；如果候选用户图允许触达不同群体，则有机会引导内容面向不同需求。可修改特征的含义同样重要，向量在单位球上移动是模型假设，不意味着实际创作者可以无成本、连续地修改所有内容属性。

一次响应预测与真实市场均衡也不同。部署评分模型后，提供者修改内容，系统观察新交互、收集新标签并重新训练，这形成连续互动，却不是EcoAgent的同一种强化学习更新。即使每次响应优化都有解，反复训练仍可能振荡，或者因参与者信息不完全而偏离预测。关于最佳响应与均衡的区别，可参见\BookSectionRef{sec:game-best-response-equilibrium}{最佳响应与均衡存在}。

原实验包含合成环境，也包含以Yelp资料为基础、再模拟内容响应和相应标签的半合成设置。真实资料提供了初始结构，不等于观察到真实创作者按式~\eqref{eq:rs-performative-response}修改内容。因此，结果支持的是在指定成本、特征可改范围和理性响应模型下的策略比较。应用到真实供给时，应另外检查响应滞后、成本异质性、竞争、进入退出及质量标签变化，不能据此宣称已经识别真实创作者策略或保证长期均衡。

\section{长期效果检验}
\label{sec:rs-long-term-evaluation}

学习到价值与证明策略改善，是两项不同工作。价值网络可以在旧策略轨迹上预测准确，新策略却访问未覆盖状态；短期点击可以提高，累计关闭、满意度或留存却朝另一个方向变化。长期检验需要固定目标人群、起点、日历窗口、交互机会与预算，并说明反馈何时成熟。

例如，假设同样的初始100位用户各观察两周，策略甲首日得到80次点击、期末有60人返回，策略乙首日得到90次点击、期末有55人返回。教学数字只说明两个指标可以方向不同，不能单靠两组总量判定因果或统计可靠性。若只比较期末仍活跃用户的平均点击，还会改变分母，把已经离开者排除在外。留存目标要求保留初始纳入人群及其后续缺失、退出口径。

模拟能够让策略在相同且已知的状态变化规则下交互，帮助检查学习机制与假设敏感性；真实试验则检验这些后果是否在实际环境出现。两种证据不能互换。历史窗口基准衡量信息条件改变后模型能否适应，近全观测选择数据衡量某些选择问题，随机曝光日志提供局部行动证据，它们都不天然包含完整的长期转移。本章讨论模拟和机制对照，第~\ref{chap:rs-evaluation}~章统一说明日志识别与实际实验协议。

\subsection{RecSim与RecSim NG：用户动态的模拟检验}
\label{sec:rs-recsim}

静态日志不能对未执行的策略重新生成用户后续反应。\term{RecSim}把这些反应写成可配置环境，使推荐策略能够提出列表、获得模拟反馈，再面对更新后的用户与目录。环境由用户状态、文档或物品、用户选择和状态转移等组件组成；配置这些组件，实际上是在说明希望检验的动态假设\scite{rs-e04}

一次交互中，环境根据当前潜在状态和目录产生策略允许看到的观测；策略根据观测选列表；选择模型决定消费对象及响应；奖励和转移模型据此更新用户、资源与目录；新的观测再进入下一步。模拟器内部可访问潜态，是为了生成过程和诊断结果，不表示推荐模型可以直接使用潜态。若给一个策略真实兴趣，另一个策略只能读日志，比较的就同时包含了信息优势。

可以在同一目录上分别建立兴趣增强和重复疲劳两个教学环境。在兴趣增强环境中，同类消费之后提高该类兴趣；在疲劳环境中，维护连续曝光计数$n_t$，令同类内容的点击概率为
\[
 p_t=\sigma(\xi_t-\kappa n_t),
 \qquad \kappa\geq0,
\]
并为中断展示后的计数恢复规定规则。若只改$\kappa$却没有说明$n_t$怎样更新与恢复，模拟并不完整。固定物品质量、初始兴趣和机会数后，可以比较点击贪心策略与轮换策略的累计收益、重复次数和离开概率；不同动态下策略顺序可能不同。

这种比较的用途是定位策略利用了哪项假设。若轮换只在很大的疲劳系数下获益，就应将该系数作为重点校准对象；若轻微改变恢复速度便反转结果，部署结论就较脆弱。环境误差还可能被策略放大：策略会不断选择模拟器预测有利的路径，从而集中访问最容易被模型误设的区域。仅让基线轨迹的平均点击接近真实数据，不足以验证这些新路径。

RecSim的可配置性不提供一个适用于所有平台的默认真实用户模型，也不能把默认环境成绩当作普遍排行榜。配置时需要记录用户与物品生成、状态可见范围、响应、转移、终止和随机源；校准时应检查分群响应和序列统计，而不只校准一个总体均值。模拟结果最终应表述为“在这些动态与信息条件下”，并与真实数据能够支持的部分分别报告。

\Needspace{100mm}
即使状态结构固定，响应参数也可能不确定。\term{RecSim NG}进一步用概率编程描述用户、提供者等实体的行为及相互依赖，把多步交互组织为可以生成和推断的概率模型。相较于只运行固定参数环境，它还允许利用轨迹对潜在状态和参数作推断；原研究展示了自动微分、哈密顿蒙特卡洛及随机期望最大化等工具的使用\scite{rs-e05}

可以将一个简化模型写成
\[
 p_\phi(s_0)
 \prod_{t=0}^{T-1}
 \pi(a_t\mid o_{\leq t},a_{<t})\,
 p_\phi(o_{t+1},r_t,s_{t+1}\mid s_t,a_t),
\]
其中$o_t$表示可见观测，$\phi$为用户和供给动态参数，初始观测机制在此从略。模型中的依赖关系决定了哪些主体先响应、哪些变量共同影响下一步。给定日志后，可以在这些依赖假设下推断未观测兴趣或疲劳，并估计参数；这不是从数据无条件恢复真实动态。

例如，重复曝光后点击减少，既可能由疲劳系数较大解释，也可能由需求已完成解释。如果模型只包含疲劳变量，后验计算再精确，也不会自动发现遗漏的需求完成机制。已知行动与观测下的良好拟合，还不能保证在未覆盖的行动序列上作出正确干预预测。

报告策略比较时，应区分参数不确定性与轨迹随机性。可以从参数后验抽取$\phi$，在每个$\phi$下再运行多个环境随机种子，计算策略差值；外层变化反映在模型内尚未确定的动态，内层变化反映同一动态下不同交互轨迹。只固定一个最优拟合参数运行许多种子，会低估前一种不确定性。两层计算仍不包括所有结构误设，兴趣、疲劳与提供者响应机制的替代模型需要单独比较。

概率模拟的价值在于把依赖、潜态和不确定性写得可检查。它不会替代行动覆盖、状态充分性或因果识别条件。真实长期对照与模拟的合理分工，是由真实反馈约束模型可接受的范围，再用模拟分析难以大量在线试探的策略与极端条件，最后对拟部署方案取得与目标相称的实际证据。

\subsection{闭环与生态效果检验}
\label{sec:rs-feedback-ecosystem-evaluation}

如果同时打开重训、兴趣转移和供给响应，再看到用户消费更集中，很难知道是哪一项机制导致。更有解释力的实验将三种机制作为独立开关，在共同初始用户、目录、机会和预算下构造对照。表~\ref{tab:rs-feedback-controls}给出基本条件；它们是实验设计，不是预先断言的结果。

\begin{table*}[htbp]
 \centering
 \small
 \begin{tabularx}{\linewidth}{@{}p{22mm}p{23mm}p{23mm}p{23mm}X@{}}
  \toprule
  条件 & 重复训练 & 兴趣转移 & 提供者响应 & 主要可检查的问题 \\
  \midrule
  静态参照 & 关闭 & 关闭 & 关闭 & 固定策略下的消费、覆盖与资源使用 \\
  数据循环 & 开启 & 关闭 & 关闭 & 策略采样与重训怎样改变交互和估计表示 \\
  用户动态 & 固定同一计划 & 开启 & 关闭 & 接触引起的兴趣、疲劳与离开怎样变化 \\
  供给动态 & 固定同一计划 & 固定同一规则 & 开启 & 机会怎样改变存续、内容与后续目录 \\
  联合动态 & 开启 & 开启 & 开启 & 多条路径之间是否存在交互作用 \\
  \bottomrule
 \end{tabularx}
 \caption{闭环与生态检验的机制对照。比较某项开关时，其余条件应相同；表中各行不是必须依次执行的算法。需要识别交互作用时，应补全相应组合，而非仅相减这些行的总结果。}
 \label{tab:rs-feedback-controls}
\end{table*}

关闭兴趣转移后，行为仍然集中，说明无需偏好变化也能产生这一现象；在其余机制一致时打开兴趣转移，差值才反映该模型中这项动态的贡献。若重训与兴趣转移同时开启后的变化不等于分别开启的变化之和，便存在交互，不能把联合变化全部归给某一项。供给进入退出改变了目录规模时，还应区分同一初始提供者的后果和新增提供者带来的构成变化。

对照不仅要共享一个整数随机种子，还应尽量对齐外生随机过程。不同策略会调用不同次数的选择与转移函数，若所有随机数来自同一顺序流，分支改变就可能打乱后续用户到达和物品质量。可以按用户到达、物品生成、响应噪声分别维护随机流，并报告多个独立重复；共同随机数用来降低比较噪声，不会让本来不可比的机制变得可比。

评价量应对应路径与主体。任务效用反映是否满足需求，行为集中度描述实际消费分布，潜在兴趣距离只在可读取或有测量模型时报告；提供者部分至少区分总效用、分布、存续以及内容数量与差异。例如，两家提供者效用从$(4,4)$变为$(9,0)$，总效用从8升到9；如果零效用者按假设退出，存续数却从2降到1。这个教学例子说明总量不能替代存续，并未规定哪一种结果应有更高优先级。

“增加供给”还需要分母。总目录变大时，覆盖物品数可能上升，覆盖比例反而下降；用户增加时，总消费上涨也可能不伴随人均任务效用上涨。为解释策略效应，应同时保留初始队列、全体动态人群和可用机会等必要口径，不在比较中临时换成更有利的指标。若策略使用更多探索或更多计算，应将探索预算、调用次数和延迟一起报告。

敏感性分析应围绕结论真正依赖的条件，而不是只扰动无关超参数。用户侧可改变强化、疲劳和恢复速度；供给侧可改变成本、机会收益曲线、退出阈值、响应滞后及跨提供者竞争；数据侧可改变覆盖、重训频率和反馈延迟。若EcoAgent的优势依赖未选者不受竞争影响，就应在逐步增加干扰的环境中检查该优势与估值偏差，而不是只在满足分解条件的环境内重复验证。

真实长期试验还有模拟对照之外的约束。个体随机化可能共享提供者与预算，一组带来的创作变化会改变另一组的目录；根据任务可能需要更大的随机化单位、市场隔离或明确的干扰分析。观察时间变长也不会自动消除季节变化、参与者退出或学习更新带来的偏差。评价时应保留预先规定的起止时间、反馈成熟窗口、纳入群体和停止规则，避免把较晚返回者误计为没有返回。

训练分布变化、策略主动改变状态与外部趋势也须分别诊断。可以记录模型版本和重训时间，追踪曝光选择与表示变化；用可测的重复频次、用户反馈和供给行为检验状态转移；再用同期对照或适当设计处理季节性与市场变化。哪些结论来自实际观测识别，哪些来自模拟机制，哪些仍需额外干预，应随结果一起说明。长期价值数值即使计算得很精确，也不能替代这条证据链。

\section{本章小结}
\label{sec:rs-long-term-summary}

当前行动影响未来，不只因为奖励会延迟到达，还因为行动会改变信息、用户状态、资源与供给。连续决策因此需要把时间范围、可见信息、转移和累计目标一起定义。长历史是输入条件，长期结果预测是估计任务，优化长期策略则要求比较改变行动之后的后果；三者不能仅凭模型读取了序列就相互替代。

价值估计与策略优化可以相互支持，但其结构条件决定了能复用什么。SlateQ在至多单消费及给定消费后的奖励、转移独立条件下分解列表价值，仍保留选择竞争与零消费；抽样策略使用的可加收益与包含概率具有不同含义。提供者增量依赖未选者累计后果不随被选身份变化，创作响应训练则依赖成本、可修改特征与行为模型。去掉这些条件，原来的计算简化不再自动成立。

是否值得部署，还取决于结论的证据范围。固定偏好的重训诊断、兴趣动态诊断和供给响应检验分别回答不同问题；模拟让机制可控，不能代替真实长期对照。把累计目标、策略变化与可识别的效果连接起来，需要进一步检查日志覆盖、随机化、反馈成熟及干扰。下一章将从这些要求出发，建立推荐评价与效果识别的共同协议。
